\documentclass[12pt]{article}
\usepackage[utf8]{inputenc}
\usepackage[T1]{fontenc}
\usepackage[margin=1in]{geometry}
\usepackage{setspace}
\usepackage{xcolor}
\usepackage[colorlinks=true, urlcolor=blue!60!black, linkcolor=blue!60!black]{hyperref}
\setlength{\parindent}{0em}
\setlength{\parskip}{0.6em}
\renewcommand{\baselinestretch}{1.0}
\begin{document}
\begin{center}
{\LARGE \textit{Trump v. United States}, 603 U.S. 593 (2024)\par}
\vspace{0.5em}
{\large Case Brief --- Two-Page Edition\par}
\end{center}

\section{Memory Jogger}
Former Presidents have absolute immunity for core powers, presumptive immunity for other official acts, and none for unofficial acts (\href{https://www.law.cornell.edu/supremecourt/text/23-939}{603 U.S.\ 593}).

\section{Detailed Case Facts}
A federal indictment charged Donald Trump with conspiring to overturn the 2020 election by pressuring state officials, organizing fraudulent electors, trying to use the Justice Department to back false fraud claims, pressuring Vice President Pence to reject electoral votes, and exploiting the January 6 attack.

\section{Procedural History}
The district court denied Trump's immunity motion, and the D.C.\ Circuit affirmed, reasoning that ``former President Trump has become citizen Trump.''

\section{Judicial Votes}
Vacated and remanded, 6--3 (5--4 on Part III-C). Roberts wrote for Thomas, Alito, Gorsuch, Kavanaugh, and Barrett (except III-C); Thomas and Barrett concurred; Sotomayor, Kagan, and Jackson dissented.

\section{Holding}
A former President is absolutely immune for acts within his ``conclusive and preclusive'' authority, at least presumptively immune for other official acts within the ``outer perimeter,'' and not immune for unofficial acts. Courts may not probe motive, and evidence of immune acts is inadmissible. Trump's Justice Department dealings were immune; other allegations were remanded.

\section{Analysis of Opinions}
Roberts relied on \textit{Youngstown}'s account of exclusive powers, Hamilton's ``energetic'' Executive, and \textit{Nixon v.\ Fitzgerald}, warning of ``an Executive Branch that cannibalizes itself.'' Thomas questioned the Special Counsel's appointment. Barrett would frame immunity as an as-applied challenge and rejected the evidentiary bar. Sotomayor called the immunity ``atextual, ahistorical, and unjustifiable,'' citing the Impeachment Judgment Clause and warning of a President ordering SEAL Team 6 to kill a rival. Jackson said the Court revived the idea that ``the King can do no wrong.''

\section{Examples: Future Applications}
\textbf{Same-Side:} corruptly ordering the Attorney General to drop an investigation is immune. \textbf{Opposite-Side:} falsifying a family business's tax returns is unofficial and not immune. \textbf{Fence-Sitter:} conditioning foreign aid on an investigation of a domestic rival.

\section{Critique}
Goldsmith, Baude, and Roisman fault the Court's reliance on structural inference and functional fears over text; Metzger calls the balancing ``faux'' and says it weakens accountability.

\section{Key Quotations}
\begin{quote}
``We conclude that under our constitutional structure of separated powers, the nature of Presidential power requires that a former President have some immunity from criminal prosecution for official acts during his tenure in office. At least with respect to the President's exercise of his core constitutional powers, this immunity must be absolute. As for his remaining official actions, he is also entitled to immunity.''
\end{quote}
\noindent --- Chief Justice Roberts, opinion of the Court.

\end{document}
